% Migrated from legacy Markdown source: 0016-retrieval-api.md (all source sections and text preserved)
% Converted to plain LaTeX for the AgentSouriva development reference.

\subsection{Retrieval API}

\commitlink{3bf0fefbf7db3f692740a7512b8f1762862624c2}{3bf0fef}

\subsubsection{Change}

Added typed evidence schemas and \texttt{POST\ /retrieve} for direct corpus search.

The route applies request bounds and the runtime relevance threshold to the same store used by the agent:

\begin{verbatim}
@app.post("/retrieve", response_model=RetrieveResponse)
def retrieve(request: RetrieveRequest) -> RetrieveResponse:
    evidence = runtime.store.search(
        request.query,
        top_k=request.top_k,
        min_score=runtime.settings.min_relevance,
    )
    return RetrieveResponse(evidence=[item.__dict__ for item in evidence])
\end{verbatim}

The response model serializes complete evidence records, letting clients inspect ranking and citations without triggering answer generation.

\subsubsection{Architecture}

The endpoint exposes the same configured FAISS retriever used by the agent and MCP tools. Results preserve the stable chunk ID, text, source, location, and measured similarity score.

\subsubsection{Design notes}

Query length and result count are bounded by request validation. The global relevance threshold is applied consistently so direct retrieval and agent decisions observe the same evidence floor.

\subsubsection{Verification}

Compiled source and tests, checked the diff, and extended the ingestion integration test to restore and query the persisted corpus. Added request validation coverage for invalid \texttt{top\_k}.

\subsubsection{Result}

Clients can inspect ranked evidence independently of answer generation.

% Optional reviewer note: \aftabnote{Write a review note here.}
